%================================================================================
% Lynceus Template Extraction (Reconstructed, EN)
% Source project : dylanyunlon/lynceus-CMD
% Template        : dylanyunlon/Desynced-Low-Communication
%                   (template_extraction_des_loc_reconstructed_EN.tex)
% Claude-1 (M001-M025): MAPPING TABLE + Section 1 (Introduction)
%                       + Section 2 (Lynceus Cost Model & Routing)
%
% The placeholder skeleton (A.., M.., I.., CI..) is inherited verbatim from the
% DES-LOC template; only the *mapping* is re-pointed to Lynceus's academic
% narrative: cost-model-driven query routing on heterogeneous GPU-CPU systems.
% No source-file / class / function names appear in the prose -- only the
% scholarly concepts they implement.
%
% Subsequent Claudes edit THIS .tex (chained patches):
%   Claude-2 -> Section 3 (Cost-Model Correctness Guarantees / Invariants)
%   Claude-3 -> Section 4 (Experimental Design)
%   Claude-4 -> Section 5 (Evaluation, RQ1-RQ6)
%   Claude-5 -> Sections 6-7 (Related Work + Conclusion)
%================================================================================

%================================================================================
% MAPPING TABLE (Section 1-2 scope; Claude-1)
%================================================================================
% Latin uppercase A-Z:
%   A = Lynceus (the cost-model-driven router)
%   B = the analytical-query cost model
%   C = single-device (GPU-only / CPU-only) execution
%   D = heterogeneous query-processing system
%   E = framework
%   F = heterogeneity-aware cost modeling
%   G = (Heterogeneity-Aware Cost Models)
%   H = static hybrid routing
%   I = cross-device data-transfer cost
%   J = relational operators
%   K = execution
%   L = data placement / residency
%   M = device-specific cost estimates
%   N = cross-device data-movement overhead
%   O = rule-based query optimization
%   P = plan-level cost estimation
%   Q = per-device minimum-cost search
%   R = per-operator device assignment
%   S = hardware devices (GPUs and CPUs)
%   T = coupling
%   U = decoupling
%   V = layers
%   W = base-table / index data
%   X = compact statistics representation
%   Y = block-wise FP8 quantization
%   Z = abstraction
%
% Greek letters:
%   alpha = admission / integration
%   beta  = variant
%   gamma = management
%   delta = component
%   epsilon = cost-model coefficients (per-row scan, seek, per-op)
%   zeta  = robustness penalty
%   eta_var = estimate (a single cost estimate)
%   theta = cardinality / selectivity estimation
%   iota  = concrete routing strategy
%   kappa = caller / query planner
%   lambda = specialization
%   mu    = decision rule
%   nu    = warm-up
%   xi    = query feature / parameter
%   omicron = unified
%   pi    = abstraction (same as Z)
%   rho_var = virtual index
%   sigma = block-wise
%   tau   = quantization
%   upsilon = data
%   phi   = pipeline / dataflow
%   chi   = configuration
%   psi_var = adoption policy
%   omega = device-load state
%
% Roman numerals:
%   I    = caching
%   II   = paged block table
%   III  = lightweight
%   IV   = physical block
%   V    = block pool
%   VI   = block identity (table, index, block-index)
%   VII  = content-addressed
%   VIII = tracking
%   IX   = allocation / eviction operations
%   X    = query routing
%   XI   = logical table identity
%   XII  = cache manager
%   XIII = reference counting
%   XIV  = channel
%   XV   = physical copy
%   XVI  = (reused) decay half-life of an exponential moving average
%   XVII = awareness
%   XVIII= adaptive frequency tuning
%   XIX  = bias-correction ratio (observed / estimated)
%   XX   = decay rate
%   XXI  = device-load / bias state
%   XXII = decoupling
%   XXIII= feedback lifecycle
%   XXIV = EMA update rule
%   XXV  = dispatch
%   XXVI = load-balancing margin
%   XXVII= queueing stall
%   XXVIII = high-frequency
%   XXIX = probing
%   XXX  = device memory (HBM)
%   XXXI = out-of-memory (OOM)
%   XXXII= distributed cost-model synchronization layer
%
% (Section 3+ mappings, e.g. invariants/contracts CI-.., are added by later
%  Claudes starting at CI -- do NOT overwrite A-XXXII above.)
%================================================================================

\documentclass{article}
\usepackage{amsmath,amssymb}

\begin{document}

%================================================================================
% Section 1: Introduction (Templated)
%================================================================================
\section{Introduction}
\label{sec:introduction}

% Paragraph 1: Background
Modern analytical data systems run on increasingly S, pairing many-core
CPUs with GPUs in a single node for improved throughput and XXX capacity.
However, routing every query to a fixed C target -- whether C GPU or C CPU --
ignores that each J has a very different cost profile on each device, and
that moving W across the device fabric incurs substantial N. Early systems
side-stepped this by hand-written H, dispatching whole queries to one device
according to a fixed rule rather than estimating cost per device.
However, modern heterogeneous workloads, e.g.\ mixed OLAP query streams, exhibit
operator-level cost asymmetry that such coarse routing cannot exploit.

% Paragraph 2: Problems with existing methods
Some approaches extend O to the heterogeneous setting by costing whole-query
plans on a single device, which poses challenges. First, they omit I,
treating data as already resident and producing optimistic estimates. Second,
charging only on-device work accumulates systematic error and provides no means
to alpha hot W into faster XXX. This makes them brittle in failure-prone or
bandwidth-constrained deployments. Third, ignoring multi-hop paths through the
device fabric destabilizes routing on multi-socket machines.

% Paragraph 3: The single-device-routing problem
H addresses some of these challenges, showing that pinning a query to the
device its bulk operator favors can beat naive C, and remain robust to adding
new S. However, it requires I to be folded into compute and never priced
separately, hiding N relative to a transfer-aware estimate. Hence, we aim to
answer the following question:

\begin{center}
\emph{Can independently pricing J K and cross-device I improve N
efficiency for heterogeneous query processing while preserving routing
correctness and robustness?}
\end{center}

% Paragraph 4: Lynceus proposal
As a result of our inquiry, we propose A, a system that drives R from a unified,
F. By estimating M on every device and routing each J to its lowest-cost target,
A reduces N by keeping hot W resident and moving data only when it pays off.

% Contributions
\paragraph{Contributions:}
\begin{enumerate}
    \item \textbf{Correct cost modeling.} We give a set of binding cost-model
    contracts (see Section~\ref{sec:guarantees}) for A: a single query's
    latency is its full critical path including I; per-query R never invents
    cross-device overlap; XI is an explicit field, not string-inferred; and
    fabric I follows multi-hop shortest paths. Our analysis indicates correct
    transfer accounting is what makes per-device R trustworthy.

    \item \textbf{N reduction.} We empirically show that data-resident routing
    and III block II reduce N by skipping redundant W movement, yielding
    end-to-end latency reductions over C on heterogeneous hardware.

    \item \textbf{Scalability across workloads.} We validate A on
    standard analytical benchmarks, demonstrating competitive latency against
    both H and C while routing each J to the cheaper device.

    \item \textbf{Hardware robustness.} Unlike fixed-rule routing, A reasons
    over the full hardware topology and prices Y, enabling it to exploit
    second-NUMA-node memory and compact statistics under tight XXX budgets.
\end{enumerate}

%================================================================================
% Section 2: Lynceus Cost Model and Routing (Templated)
%================================================================================
\section{A}
\label{sec:lynceus}

% Paragraph 1: Cost-asymmetry theory
We start by characterizing the relation between where W resides and the cost of
running J on each device, and how this can be leveraged to lower N. Consider the
per-device M as a sum of itemized terms:
\begin{equation}
\text{cost}(J, S) = c_{\text{io}} + c_{\text{compute}} + c_{\text{transfer}} + c_{\text{index}} + c_{\text{sort}}.
\end{equation}

For C routing, the choice is correct only if $c_{\text{transfer}}$ is counted:
a GPU may win on $c_{\text{compute}}$ yet lose once the PCIe load into XXX is
added. Conversely, if W is already resident the transfer term vanishes and the
GPU path dominates on far more J.

% Paragraph 2: Itemized device cost
A useful summary is the per-device cost shape: a CPU prices a scan by
$c_{\text{io}} \propto$ rows $\times$ ε, while a GPU amortizes the same
scan over massive parallelism but pays a fixed I to bring the data on-device.
We use these ε to compare devices on a common microsecond scale. For
typical coefficients, a sequential CPU scan and a transfer-plus-parallel GPU
scan cross over at a workload-dependent cardinality, which the cost model
locates per query.

% Paragraph 3: Transfer cost over the fabric
Because W may sit on a remote node, I is not a single edge but the best path
through the device fabric. On a dual-socket box the second NUMA node may have no
direct GPU link, so data reaches a GPU via a multi-hop route. The fabric cost is
therefore the minimum-cost path, additive over hops, with disconnected cuts
priced as unreachable.

% Paragraph 4: Routing decision and half-life-style adaptation
From the above, the resident location of W should inform R. For example, if hot
W already occupies XXX, the transfer term drops out and routing J there is
cheap. Where the static M is biased, an online policy corrects it: tracking a
per-device XIX with decay rate XX and adapting XXV accordingly.

%-----------------------------------------------------------------------------
\subsection{The A Router}
\label{sec:lynceus_router}

% Router / strategy description
A routes each J by estimating M on every device $S$ and selecting the
lowest-cost target, optionally pricing ρ so a J that hits an index pays a
traversal rather than a full scan. Setting the policy to pure minimum-cost
yields cost-model routing; adding a ζ that prefers the robust device when the
margin is thin yields a penalty-aware variant.

\paragraph{Illustrative Example} Figure~\ref{fig:toy} illustrates a workload
where A and a transfer-aware H route correctly under heterogeneous device costs,
while fixed C routing -- which ignores I -- mispredicts and overloads one device.

\end{document}
